\noindent\textbf{Elección del procedimiento.} La hidrostática utiliza altura vertical, aunque la escala mida a lo largo del tubo: $h=L\sin\theta$. El aumento de longitud para un mismo desnivel es $L/h=1/\sin\theta$.

\begin{enumerate}

\item Proyectar la lectura sobre la vertical.

\[\fitmath{h=L\cdot \sin\left(\theta\right)}\]

\[\fitmath{h=(0.12)\cdot \sin\left((0.261799)\right)\approx 0.03105824\,\mathrm{m}}\]

\item Comparar la lectura oblicua con su componente vertical.

\[\fitmath{\mathrm{amp}=\frac{1}{\sin\left(\theta\right)}}\]

\[\fitmath{\mathrm{amp}=\frac{1}{\sin\left((0.261799)\right)}\approx 3.863709\,}\]

\item Aplicar el balance hidrostático al desnivel obtenido.

\[\fitmath{\mathrm{dp}=\mathrm{rhom}\cdot g\cdot h}\]

\[\fitmath{\mathrm{dp}=(790)\cdot (9.81)\cdot (0.03105824046)\approx 240.6983\,\mathrm{Pa}}\]

\end{enumerate}
